\begin{pastexam}[label={px:ncu-114-dsa-14}]{中央資工所 114 學年度 第 14 題}
（複選；每題 5 分，全部答對才給分，答錯倒扣 1 分）

Which of the following statements about the Maximum Flow problem are incorrect?
\begin{enumerate}[label=(\alph*)]
\item In a flow network, source node has no incoming flow and terminal (sink) node has no outgoing flow.
\item Each edge has forward flow and backflow, and the two flows must always be inverses of each other.
\item Ford-Fulkerson algorithm finds the maximum flow from source to terminal with time complexity
$O(Ef)$, where $E$ represents the number of edges and $f$ is the maximum flow of the final graph.
\item If we look at all the possible cuts in the graph, and find the largest capacity of those cuts,
then that value is the value of the maximum flow for that network.
\item If a network has antiparallel edges (i.e., the edges connecting two nodes back and forth), we can
convert it to an equivalent one with no antiparallel edges.
\end{enumerate}
\end{pastexam}

\begin{pastsoln}
答案：(d)。 題目問的是\textbf{錯誤}的敘述。

(a)(b) 兩句是某些教材的定義或約定，對錯取決於採用哪一套定義，先說明本題的讀法。 (a) 是 Kleinberg 與
Tardos 等教材對流網路的假設：沒有邊進入源點、沒有邊離開匯點。 (b) 說的是「淨流」寫法（CLRS 第二版）
的斜對稱性 $f(u, v) = -f(v, u)$：$u \to v$ 送 $3$ 單位，就記成 $v \to u$ 送 $-3$ 單位。 照這些定義，
兩句都是對的。 資結書第~26.1~節的定義比較一般：允許有邊指進源點，流量寫成 $0 \le f(u, v) \le c(u, v)$
而不用斜對稱；在那套定義下，(a)(b) 不是定義的一部分。 本題照出題的語氣，把 (a)(b) 當成正確的敘述。

\begin{itemize}
\item (a) 正確（見上）。 就算網路有邊指進 $s$，也一定有一個最大流在這些邊上流量為 $0$，所以這個約定
不影響最大流的值。
\item (b) 正確（見上）。 在剩餘網路裡看也是同一件事：邊 $(u, v)$ 上流了 $f$，就多一條反向的剩餘邊
$(v, u)$，容量 $f$，代表可以把這些流「退回去」（資結書第~26.2~節）。
\item (c) 正確。 容量是整數時，每次增廣至少增加 $1$ 單位，最多 $|f^*|$ 次；每次用 BFS 或 DFS 在剩餘
網路找路徑 $O(E)$，總共 $O(E\,|f^*|)$（資結書第~26.3~節）。 （容量是無理數時 Ford--Fulkerson 可能
不會停，題目的 $O(Ef)$ 隱含整數容量。）
\item (d) 錯誤。 最大流最小割定理（資結書第~26.5~節）：最大流的值等於所有割裡\textbf{最小}的容量，
不是最大的。 任何流的值都 $\le$ 任何割的容量，所以最大流只會碰到最小的那個割。 反例：$s \to a$ 容量
$1$、$a \to t$ 容量 $5$，割 $(\{s\}, \{a, t\})$ 容量 $1$、割 $(\{s, a\}, \{t\})$ 容量 $5$，最大流是 $1$，
不是 $5$。
\item (e) 正確。 有 $u \to v$ 和 $v \to u$ 時，挑一條，例如 $(v, u)$，加一個新頂點 $x$，換成
$v \to x$、$x \to u$，容量都是原來的 $c(v, u)$（資結書第~26.1.2~節）。 新網路的最大流值和原網路一樣。
\end{itemize}
\end{pastsoln}
